% title:          Diagonal products versus eigenvalue products
% area:           bilinear-forms, eigenvalues, trace
% methods:        trace-argument, eigenvalue-arguments
% difficulty:
% difficulty-ai:  easy
% status:
% review:         none
% --- history: all optional, leave EMPTY if unknown ---
% origin:         imc
% origin-date:    2014-08-01
% origin-number:  Day 2, Problem 2
% also-in:
% taken-from:
% references:     IMC 2014 (21st IMC, Blagoevgrad), Day 2, Problem 2, proposed by Martin Niepel (Comenius University, Bratislava) - https://www.imc-math.org.uk/?year=2014&item=problems ; questions https://www.imc-math.org.uk/imc2014/IMC2014-day2-questions.pdf ; official solution https://www.imc-math.org.uk/imc2014/IMC2014-day2-solutions.pdf . Mathematical background: equivalent to sum a_ii^2 <= sum lambda_i^2, a consequence of Schur's 1923 diagonal/eigenvalue inequalities (N. Higham, 'Eigenvalue Inequalities for Hermitian Matrices', 2021 - https://nhigham.com/2021/03/09/eigenvalue-inequalities-for-hermitian-matrices/)
% history:        ai
\begin{problem}
Let \( A = (a_{ij})_{i,j=1}^{n} \) be a symmetric \( n \times n \) matrix with real entries, and let \( \lambda_{1}, \lambda_{2}, \ldots, \lambda_{n} \) denote its eigenvalues. Show that
\[
\sum_{1 \leq i < j \leq n} a_{ii} a_{jj} \geq \sum_{1 \leq i < j \leq n} \lambda_{i} \lambda_{j},
\]
and determine all matrices for which equality holds.
\end{problem}
\begin{solution}
Eigenvalues of a real symmetric matrix are real, hence the inequality is well-defined.
\\
The trace of a matrix equals the sum of its eigenvalues. For matrix \( A \),
\[
\sum_{i=1}^{n} a_{ii} = \sum_{i=1}^{n} \lambda_{i}.
\]
Squaring both sides, we obtain:
\[
\left( \sum_{i=1}^{n} a_{ii} \right)^2 = \left( \sum_{i=1}^{n} \lambda_{i} \right)^2.
\]
Expanding both sides gives:
\[
\sum_{i=1}^{n} a_{ii}^2 + 2 \sum_{1 \leq i < j \leq n} a_{ii} a_{jj} = \sum_{i=1}^{n} \lambda_{i}^2 + 2 \sum_{1 \leq i < j \leq n} \lambda_{i} \lambda_{j}.
\]
It therefore suffices to show the inequality:
\[
\sum_{i=1}^{n} a_{ii}^2 \leq \sum_{i=1}^{n} \lambda_{i}^2.
\]
The matrix \( A^2 \), which equals \( A^T A \) for symmetric \( A \), has eigenvalues \( \lambda_{1}^2, \lambda_{2}^2, \ldots, \lambda_{n}^2 \) for the same reason as the previous problem. The trace of \( A^T A \) is the square of the Frobenius norm of \( A \):
\[
\operatorname{tr}\left(A^T A\right) = \sum_{i,j=1}^{n} a_{ij}^2 = \operatorname{tr}\left(A^2\right) = \sum_{i=1}^{n} \lambda_{i}^2.
\]
Obviously \( \sum_{i=1}^{n} a_{ii}^2 \leq \sum_{i,j=1}^{n} a_{ij}^2 \), the inequality \( \sum_{i=1}^{n} a_{ii}^2 \leq \sum_{i=1}^{n} \lambda_{i}^2 \) follows. One sees then that equality holds if and only if $\sum_{i=1}^{n} a_{ii}^2 = \sum_{i,j=1}^{n} a_{ij}^2$ that is if and only if all off-diagonal entries of \( A \) are zero, i.e., \( A \) is diagonal.
\begin{remark}
The same result holds for Hermitian matrices as for Hermitian matrices, diagonal entries and eigenvalues are also real. 
\end{remark}
\end{solution}
